\documentclass{report}

\begin{document}
	\section*{Federico's notes}
	\subsection*{Progress of the paper}
	\begin{enumerate}
		\item The introduction should stress that 1) What is the realization of an augmented simplicial set? A natural way to realize it is to define it as the homotopy cofiber of the map induced on realization by the augmentation. This produces a space, but without cup products (one gets disjoint unions of suspensions of spaces). Endowing them with cup-i products is new, I believe. 2) There are other interesting categories that may not have a realization into spaces. 3) Stress that our construction is completely independent of the Alexander-Whitney diagonal (which is quite difficult to recover) and the Surjection operad approach.
		\item Establish the sign that makes the operations commute up to suspension (this is probably another section).
		\item Rewrite Section 2 (small changes)
		\item Rewrite Section 3 adding what we know about the stable and the Tate resolution.
		\item In Section 5 and 6, talk about Tate resolutions and possibly (only if it is worthy) about products.
		\item Sections 6.2, 6.3 and 6.4 should be a separate section.
		\item Is it good to have a section on Kan spectra? I like the idea, but not sure whether it is worth it.
		\begin{itemize}
			\item Good: The operations are specially well-suited for Kan spectra: With the current notation we only need to replace the finite joins by infinite joins et voilá.
			\item Bad: Nobody uses Kan spectra.
			\item Good: I like Kan spectra.
			\item Bad: The longer the section, the worse: We don't want this to be an introduction to Kan spectra.
		\end{itemize}
	\end{enumerate}\title{Federico's notes}
	\subsection*{Curiosities}
	\begin{enumerate}
		\item Try to prove that the operation in degree $0$ induces a product in cohomology.
		\item Prove that the operation that computes the Bockstein (the one induced by $\Psi^\vee(e_r^\vee)$) does indeed compute the Bockstein.
		\item Prove a Cartan formula for the join product.
		\item Try to prove the Adem relations as Steenrod does in `Cohomology operations''. In that respect, notice that the map $\alpha$ is not only $C_r$-equivariant, but also $\Sigma_r$-equivariant. What can be said about $\psi$? Both the source and target have an action of the symmetric group, but the construction seems to be equivariant only for the action of the cyclic group.
	\end{enumerate}
	\subsection*{Further questions}
	\begin{enumerate}
		\item Rephrase the maps $\phi, \psi$ conjugating by $\Lambda$.
		\item We seem to be proving that the Eilenberg-Zilber simplicial operad in arity $r$ is isomorphic to the chain complex of the unreduced join product of $\partial \Delta^{r-1}$ with itself. Is this true? The join product on the latter seems to play an important role in defining $\psi$, can it be related to something in the operadic structure of the Eilenberg-Zilber?
		\item Our construction of $\psi$ requires that $r$ is prime (I believe the construction of $\phi$ and $\alpha$ work for any odd $r$). Is it possible to extend it to other odd $r$'s?
		\item What can be said about even $r$'s? In particular, try to find the signs for $r=2$ (this may involve working with twisted coefficients).
		\item Rethink the construction of $\psi$, trying to relate it as much as possible to the join product and its adjoints.
		\item The fact that $\psi$ factors through the reduced join gives some properties to the cup-(p,i) products. Are they related to Anibal's properties that characterize Steenrod's cup-(2,i) products? 
	\end{enumerate}
\end{document}